\subsection{Definite Iteration: The \texttt{for} Loop}

A \texttt{for} loop is Python's main syntax for definite iteration. Definite does not mean that the number of iterations must be written as a literal number. It means that the loop receives its values from a traversal source: a string, list, tuple, range, or any other iterable object.

A \texttt{while} loop repeats while a condition remains true. A \texttt{for} loop repeats by requesting successive values from an object that knows how to provide traversal.

\begin{quote}
A \texttt{for} loop does not count by itself. It asks an iterable object for values until that iterable's iterator reports that no more values exist.
\end{quote}

The ordinary syntax is:

\begin{verbatim}
for name in iterable:
    loop_body
\end{verbatim}

The name on the left is rebound for each produced value. The object on the right supplies the traversal stream.

\subsubsection{Abstracting Sequential Traversal: Structural Syntax over Strings, Lists, Tuples, and Range Objects}

The strength of the \texttt{for} loop is that the same syntax works over many different object types. The loop body does not need to know whether the values came from a string, a list, a tuple, or a \texttt{range}. It only receives one value at a time.

A string traversal produces characters:

\begin{verbatim}
text = "Ni!"

print(f"type(text): {type(text)}")
print(f"dir(text) sample: {dir(text)[:10]}")
print(f"len(text): {len(text)}")

for ch in text:
    print(f"character: {ch!r}")
\end{verbatim}

Possible output:

\begin{verbatim}
type(text): <class 'str'>
dir(text) sample: ['__add__', '__class__', '__contains__', '__delattr__',
'__dir__', '__doc__', '__eq__', '__format__', '__ge__', '__getattribute__']
len(text): 3
character: 'N'
character: 'i'
character: '!'
\end{verbatim}

The string object supports traversal because it is iterable. Its \texttt{\_\_iter\_\_()} method can produce an iterator over its characters. Its \texttt{\_\_len\_\_()} method gives its length, and \texttt{\_\_getitem\_\_()} supports indexed access.

A list traversal produces stored elements:

\begin{verbatim}
quotes = [
    "No soup for you!",
    "Giddy up!",
    "D'oh!",
]

print(f"type(quotes): {type(quotes)}")
print(f"dir(quotes) sample: {dir(quotes)[:10]}")
print(f"len(quotes): {len(quotes)}")
print(f"has append: {'append' in dir(quotes)}")

for quote in quotes:
    print(f"quote: {quote}")
\end{verbatim}

Possible output:

\begin{verbatim}
type(quotes): <class 'list'>
dir(quotes) sample: ['__add__', '__class__', '__class_getitem__',
'__contains__', '__delattr__', '__delitem__', '__dir__', '__doc__',
'__eq__', '__format__']
len(quotes): 3
has append: True
quote: No soup for you!
quote: Giddy up!
quote: D'oh!
\end{verbatim}

A list is a mutable container. Its \texttt{append()} method can add elements, its \texttt{\_\_getitem\_\_()} method supports indexing, and its \texttt{\_\_iter\_\_()} method can produce a traversal object.

A tuple traversal looks similar, but the tuple itself is immutable:

\begin{verbatim}
facts = (
    "Chuck Norris counted to infinity.",
    "Twice.",
    42,
)

print(f"type(facts): {type(facts)}")
print(f"dir(facts) sample: {dir(facts)[:10]}")
print(f"len(facts): {len(facts)}")
print(f"has append: {'append' in dir(facts)}")

for fact in facts:
    print(f"fact: {fact!r}")
\end{verbatim}

Possible output:

\begin{verbatim}
type(facts): <class 'tuple'>
dir(facts) sample: ['__add__', '__class__', '__class_getitem__',
'__contains__', '__delattr__', '__dir__', '__doc__', '__eq__',
'__format__', '__ge__']
len(facts): 3
has append: False
fact: 'Chuck Norris counted to infinity.'
fact: 'Twice.'
fact: 42
\end{verbatim}

The tuple can be traversed, but it cannot be structurally modified like a list. Traversal and mutability are separate properties. An object may be iterable without being mutable.

A \texttt{range} traversal produces computed integers:

\begin{verbatim}
steps = range(1, 4)

print(f"type(steps): {type(steps)}")
print(f"dir(steps) sample: {dir(steps)[:10]}")
print(f"steps.start: {steps.start}")
print(f"steps.stop: {steps.stop}")
print(f"steps.step: {steps.step}")

for step in steps:
    print(f"step {step}: And now for something completely different.")
\end{verbatim}

Possible output:

\begin{verbatim}
type(steps): <class 'range'>
dir(steps) sample: ['__bool__', '__class__', '__contains__',
'__delattr__', '__dir__', '__doc__', '__eq__', '__format__',
'__ge__', '__getattribute__']
steps.start: 1
steps.stop: 4
steps.step: 1
step 1: And now for something completely different.
step 2: And now for something completely different.
step 3: And now for something completely different.
\end{verbatim}

The \texttt{range} object does not store the produced integers as a list. It stores the arithmetic rule. The \texttt{for} loop does not care. It only receives the next produced value.

The same loop shape therefore covers several storage models:

\begin{verbatim}
for x in "abc":
    print(x)

for x in ["a", "b", "c"]:
    print(x)

for x in ("a", "b", "c"):
    print(x)

for x in range(3):
    print(x)
\end{verbatim}

The surface syntax is the same. The traversal source changes.

The loop variable is rebound on each iteration:

\begin{verbatim}
items = ["soup", "pretzels", "muffin"]

for item in items:
    print(f"inside loop item: {item}")

print(f"after loop item: {item}")
\end{verbatim}

Possible output:

\begin{verbatim}
inside loop item: soup
inside loop item: pretzels
inside loop item: muffin
after loop item: muffin
\end{verbatim}

The name \texttt{item} is not a new private slot created only for the loop body. It is a normal Python name in the surrounding scope. After the loop finishes, it remains bound to the last value produced by the traversal.

This is important for programmers coming from C-style loop headers. Python's \texttt{for} loop is not primarily:

\begin{verbatim}
initialize counter
test counter
increment counter
\end{verbatim}

It is primarily:

\begin{verbatim}
obtain iterator from iterable
ask iterator for next value
bind value to loop name
execute body
repeat until no more values exist
\end{verbatim}

When numeric positions are needed, \texttt{range()} supplies them:

\begin{verbatim}
names = ["Jerry", "Elaine", "Kramer"]

for i in range(len(names)):
    print(f"{i:>2} | {names[i]}")
\end{verbatim}

Possible output:

\begin{verbatim}
 0 | Jerry
 1 | Elaine
 2 | Kramer
\end{verbatim}

This works, but if both index and value are needed, Python commonly uses \texttt{enumerate()}, which will be examined later.

For now, the basic rule is:

\begin{quote}
A \texttt{for} loop abstracts traversal. It hides the repeated calls needed to fetch successive values from the source object.
\end{quote}

\subsubsection{Under the Hood: How CPython Implicitly Calls \texttt{iter()} and \texttt{next()} Until \texttt{StopIteration}}

The \texttt{for} loop is built on the iterator protocol. At the language level, this means two operations matter:

\begin{itemize}
    \item \texttt{iter(obj)} asks an iterable object for an iterator.
    \item \texttt{next(iterator)} asks that iterator for the next value.
\end{itemize}

When the iterator has no more values, it raises \texttt{StopIteration}. The \texttt{for} loop catches that signal internally and exits normally.

A normal \texttt{for} loop:

\begin{verbatim}
items = ["soup", "pretzels", "muffin"]

for item in items:
    print(f"item: {item}")

print("loop finished")
\end{verbatim}

Possible output:

\begin{verbatim}
item: soup
item: pretzels
item: muffin
loop finished
\end{verbatim}

A manual version using \texttt{iter()} and \texttt{next()}:

\begin{verbatim}
items = ["soup", "pretzels", "muffin"]

it = iter(items)

while True:
    try:
        item = next(it)
    except StopIteration:
        break

    print(f"item: {item}")

print("loop finished")
\end{verbatim}

Possible output:

\begin{verbatim}
item: soup
item: pretzels
item: muffin
loop finished
\end{verbatim}

The two versions have the same traversal result. The \texttt{for} loop is the clean syntax. The manual version exposes the mechanism.

The iterator object is separate from the list:

\begin{verbatim}
items = ["soup", "pretzels", "muffin"]
it = iter(items)

print(f"type(items): {type(items)}")
print(f"type(it): {type(it)}")
print(f"items is it: {items is it}")

print(f"has __iter__ on items: {'__iter__' in dir(items)}")
print(f"has __next__ on items: {'__next__' in dir(items)}")

print(f"has __iter__ on it: {'__iter__' in dir(it)}")
print(f"has __next__ on it: {'__next__' in dir(it)}")
\end{verbatim}

Possible output:

\begin{verbatim}
type(items): <class 'list'>
type(it): <class 'list_iterator'>
items is it: False
has __iter__ on items: True
has __next__ on items: False
has __iter__ on it: True
has __next__ on it: True
\end{verbatim}

The list is iterable because it has \texttt{\_\_iter\_\_()}. The list iterator is an iterator because it has \texttt{\_\_next\_\_()}. The iterator stores traversal position. The list stores the elements.

Calling \texttt{next()} advances the iterator:

\begin{verbatim}
items = ["soup", "pretzels", "muffin"]
it = iter(items)

print(f"next(it): {next(it)}")
print(f"next(it): {next(it)}")
print(f"next(it): {next(it)}")
\end{verbatim}

Possible output:

\begin{verbatim}
next(it): soup
next(it): pretzels
next(it): muffin
\end{verbatim}

One more call raises \texttt{StopIteration}:

\begin{verbatim}
items = ["soup"]
it = iter(items)

print(f"next(it): {next(it)}")

try:
    print(next(it))
except StopIteration as err:
    print(f"type(err): {type(err)}")
    print(f"message: {err}")
    print(f"dir(err) sample: {dir(err)[:8]}")
\end{verbatim}

Possible output:

\begin{verbatim}
next(it): soup
type(err): <class 'StopIteration'>
message:
dir(err) sample: ['__cause__', '__class__', '__context__', '__delattr__',
'__dict__', '__dir__', '__doc__', '__eq__']
\end{verbatim}

The \texttt{StopIteration} object is an exception object. Its role here is not to report a program failure to the user. It is the internal signal that traversal is complete. A \texttt{for} loop handles that signal automatically.

The rough expansion of a \texttt{for} loop is therefore:

\begin{verbatim}
_iterator = iter(iterable)

while True:
    try:
        target_name = next(_iterator)
    except StopIteration:
        break

    loop_body
\end{verbatim}

CPython implements this through bytecode instructions rather than literally rewriting the source into this code. Conceptually, however, the protocol is the same: get an iterator, request values, stop when \texttt{StopIteration} appears.

A \texttt{range} object follows the same mechanism:

\begin{verbatim}
nums = range(1, 4)
it = iter(nums)

print(f"type(nums): {type(nums)}")
print(f"type(it): {type(it)}")

print(f"next(it): {next(it)}")
print(f"next(it): {next(it)}")
print(f"next(it): {next(it)}")
\end{verbatim}

Possible output:

\begin{verbatim}
type(nums): <class 'range'>
type(it): <class 'range_iterator'>
next(it): 1
next(it): 2
next(it): 3
\end{verbatim}

The source object is a \texttt{range}. The traversal object is a \texttt{range\_iterator}. The range remains reusable; the iterator advances.

A tuple also produces a separate iterator:

\begin{verbatim}
data = ("red", "green", "blue")
it = iter(data)

print(f"type(data): {type(data)}")
print(f"type(it): {type(it)}")

print(f"next(it): {next(it)}")
print(f"next(it): {next(it)}")
\end{verbatim}

Possible output:

\begin{verbatim}
type(data): <class 'tuple'>
type(it): <class 'tuple_iterator'>
next(it): red
next(it): green
\end{verbatim}

For strings, the exact iterator type name may differ between Python versions and builds, but the protocol remains the same:

\begin{verbatim}
text = "abc"
it = iter(text)

print(f"type(text): {type(text)}")
print(f"type(it): {type(it)}")

print(f"has __next__: {'__next__' in dir(it)}")

print(f"next(it): {next(it)}")
print(f"next(it): {next(it)}")
print(f"next(it): {next(it)}")
\end{verbatim}

Possible output:

\begin{verbatim}
type(text): <class 'str'>
type(it): <class 'str_ascii_iterator'>
has __next__: True
next(it): a
next(it): b
next(it): c
\end{verbatim}

On some versions, the displayed iterator class name may be different, such as \texttt{str\_iterator}. The important property is not the exact printed class name. The important property is that the object returned by \texttt{iter(text)} has \texttt{\_\_next\_\_()} and advances one character at a time.

The \texttt{for} loop hides this repeated protocol call:

\begin{verbatim}
text = "abc"

for ch in text:
    print(f"ch: {ch}")
\end{verbatim}

Possible output:

\begin{verbatim}
ch: a
ch: b
ch: c
\end{verbatim}

The loop performs the repeated \texttt{next()} calls and exits cleanly when the iterator raises \texttt{StopIteration}.

This gives the exact mental model needed for later sections:

\begin{quote}
An iterable can produce an iterator. An iterator produces values one at a time. A \texttt{for} loop connects these two operations and treats \texttt{StopIteration} as normal loop completion.
\end{quote}